\documentclass[10pt,a4paper]{amsart}
\usepackage[margin=19mm]{geometry}
\usepackage{amsmath,amssymb,mathtools}
\usepackage[scr=rsfs,cal=euler]{mathalfa}
\usepackage{xfrac}
\usepackage[unicode,hidelinks,bookmarks=false]{hyperref}
\newcommand{\ie}{i.e.\ }
\newcommand{\cf}{cf.\ }
\newcommand{\resp}{respectively\ }
\newcommand{\Subsection}{\subsection}
\providecommand{\nobreakdash}{\nobreak\hbox{-}}
\setlength{\emergencystretch}{2em}
\multlinegap0pt
\usepackage{mathtools}
\usepackage{amssymb}
\usepackage{bm}
\usepackage{xcolor}
\newtheorem{proposition}{Proposition}
\newtheorem{theorem}{Theorem}
\title{\textbf{Parametric ROC Analysis and Optimal Cutoff Selection under Scale Mixtures of Skew-Normal Distributions: A Decision-Theoretic Framework with Asymptotic Inference}}
\author{Renato de Paula, Helena Mouri\~no, Tiago Dias Domingues}
\date{Source: arXiv:2605.07829v1 (CC BY 4.0)}
\begin{document}
\maketitle

\noindent\emph{Source and licence.} \textbf{Parametric ROC Analysis and Optimal Cutoff Selection under Scale Mixtures of Skew-Normal Distributions: A Decision-Theoretic Framework with Asymptotic Inference}. Version arXiv:2605.07829v1 (CC BY 4.0), distributed by arXiv under the Creative Commons Attribution 4.0 International licence, http://creativecommons.org/licenses/by/4.0/, as recorded in the arXiv licence metadata for this exact version and shown on the abstract page https://arxiv.org/abs/2605.07829v1. Full licence notice: \texttt{licenses/arxiv-2605.07829-CC-BY-4.0.txt}.\par\smallskip

\noindent\textit{Editorial scope.} Counted unit: the article's theorem thm:asymptotics (source0.tex lines 1522--1570). The article's complete original proof of it is delivered with it (source0.tex lines 1572--1685, one \texttt{proof} environment of 114 lines). No mathematics is written, repaired or re-derived: the statement and the proof are the article's own lines, in the article's own notation, and no retained line is substituted or re-flowed.  Retained uncounted as the source-local context the proof needs: lines 1046--1082; lines 1481--1486.  Nothing was omitted from any retained span, so the package carries no omission record.  No retained line contains a citation, so the wrapper carries no bibliography and no bibliography entry.  No retained line contains a figure environment, an image-inclusion command or drawing code, so no image is delivered and the package declares no asset dependency.  Editorial formatting only, in addition: an AMS \texttt{amsart} preamble that restates the article's own declarations and macros used by the retained lines, namely the environments \texttt{proposition}, \texttt{theorem} on separate counters, together with the standard packages those lines need.  The retained environments print in this document as Proposition 1, Theorem 1 in order of appearance, whereas the article prints the same items with the numbering of its own full document.  The complete native source of the article as distributed in the arXiv e-print for this exact version is bundled unchanged as \texttt{sources/200113/source0.tex}.
\begin{proposition}[\textbf{Existence, uniqueness and global optimality under MLR}]
\label{prop:optimal-cutoff}
Assume that:
\begin{enumerate}
    \item the densities $f_0(\cdot\,;\theta_0)$ and
          $f_1(\cdot\,;\theta_1)$ are continuous and strictly
          positive on $(a,b)$;

    \item the likelihood ratio
          \[
          \Lambda(c;\theta)
          =
          \frac{f_1(c;\theta_1)}{f_0(c;\theta_0)}
          \]
          is continuous and strictly increasing on $(a,b)$;

    \item the boundary conditions
          \[
          \lim_{c\,\downarrow\,a}\Lambda(c;\theta)
          <
          \frac{\lambda_0\pi_0}{\lambda_1\pi_1}
          \qquad\text{and}\qquad
          \lim_{c\,\uparrow\,b}\Lambda(c;\theta)
          >
          \frac{\lambda_0\pi_0}{\lambda_1\pi_1}
          \]
          hold.
\end{enumerate}
Then there exists a unique $c^*(\theta)\in(a,b)$ satisfying
\[
\Lambda(c^*(\theta);\theta)
=
\frac{\lambda_0\pi_0}{\lambda_1\pi_1}.
\]
Moreover, $c^*(\theta)$ is the unique global minimiser of
$R(c;\theta)$ over $[a,b]$.
\end{proposition}

\begin{equation}
\label{eq:theta-an}
\sqrt{n}\bigl(\hat\theta-\theta^*\bigr)
\stackrel{d}{\longrightarrow}
\mathcal N(0,\Sigma),
\end{equation}

\begin{theorem}[\textbf{Consistency and asymptotic normality}]
\label{thm:asymptotics}
Assume the conditions of Proposition~\ref{prop:optimal-cutoff}. Suppose further that:
\begin{enumerate}
\item the map $(c,\theta)\mapsto\varphi(c,\theta)$ is continuously differentiable in a neighbourhood of $(c^*,\theta^*)$, where $c^*=c^*(\theta^*)$;
\item
\[
\partial_c\varphi(c^*,\theta^*)\neq 0;
\]
\item the estimator $\hat\theta$ satisfies
\[
\sqrt{n}(\hat\theta-\theta^*)\stackrel{d}{\longrightarrow}\mathcal N(0,\Sigma).
\]
\end{enumerate}
Then:
\begin{enumerate}
\item $\hat c\xrightarrow{P}c^*$;
\item
\[
\sqrt{n}(\hat c-c^*)\stackrel{d}{\longrightarrow}\mathcal N(0,V),
\]
where
\begin{equation}
\label{eq:V-general}
V
=
\nabla_\theta c^*(\theta^*)^\top\Sigma\,\nabla_\theta c^*(\theta^*),
\end{equation}
and
\begin{equation}
\label{eq:gradient-cstar}
\nabla_\theta c^*(\theta^*)
= - \frac{\nabla_\theta\varphi(c^*,\theta^*)}{\partial_c\varphi(c^*,\theta^*)}.
\end{equation}
Equivalently,
\begin{equation}
\label{eq:variance}
V
=
\frac{
\bigl[\nabla_\theta\varphi(c^*,\theta^*)\bigr]^\top
\Sigma
\bigl[\nabla_\theta\varphi(c^*,\theta^*)\bigr]
}{
\bigl[\partial_c\varphi(c^*,\theta^*)\bigr]^2
}.
\end{equation}
\end{enumerate}
\end{theorem}

\begin{proof}
Since $\varphi(c^*,\theta^*)=0$ and $\partial_c\varphi(c^*,\theta^*)\neq 0$, the implicit function theorem guarantees the existence of a neighbourhood $\mathcal U$ of $\theta^*$ and a unique continuously differentiable function
\[
\theta\mapsto c^*(\theta)
\]
such that
\[
\varphi(c^*(\theta),\theta)=0,
\qquad
\theta\in\mathcal U,
\]
with $c^*(\theta^*)=c^*$.

\noindent
Because $\hat\theta\xrightarrow{P}\theta^*$ and $c^*(\cdot)$ is continuous at $\theta^*$, the continuous mapping theorem yields
\[
\hat c=c^*(\hat\theta)\xrightarrow{P}c^*(\theta^*)=c^*.
\]
This proves consistency.

For asymptotic normality, apply a first-order Taylor expansion of the differentiable map $\theta \mapsto c^*(\theta)$ around $\theta^*$:
\[
c^*(\hat\theta)-c^*(\theta^*)
=
\nabla_\theta c^*(\theta^*)^\top(\hat\theta-\theta^*)
+
\ell (\hat\theta),
\]
where the remainder term $\ell(\hat\theta)$ satisfies
\[
\frac{\ell(\hat\theta)}{\|\hat\theta-\theta^*\|} \xrightarrow{P} 0.
\]

\noindent
Since $\hat\theta$ is $\sqrt{n}$-consistent, i.e.,
\[
\hat\theta - \theta^* = O_p(n^{-1/2}),
\]
it follows that
\[
\ell(\hat\theta)
=
\|\hat\theta-\theta^*\|
\cdot
\frac{\ell(\hat\theta)}{\|\hat\theta-\theta^*\|}
=
O_p(n^{-1/2}) \cdot o_p(1)
=
o_p(n^{-1/2}).
\]

\noindent
Therefore,
\[
c^*(\hat\theta)-c^*(\theta^*)
=
\nabla_\theta c^*(\theta^*)^\top(\hat\theta-\theta^*)
+
o_p(n^{-1/2}).
\]

\noindent
Multiplying both sides by $\sqrt{n}$ yields
\[
\sqrt{n}(\hat c-c^*)
=
\nabla_\theta c^*(\theta^*)^\top \sqrt{n}(\hat\theta-\theta^*)
+
o_p(1).
\]
By~\eqref{eq:theta-an} and the multivariate delta method,
\[
\sqrt n(\hat c-c^*)
\stackrel{d}{\longrightarrow}
\mathcal N\!\left(
0,
\nabla_\theta c^*(\theta^*)^\top\Sigma\,\nabla_\theta c^*(\theta^*)
\right),
\]
which proves~\eqref{eq:V-general}.

\noindent
It remains to compute $\nabla_\theta c^*(\theta^*)$. Differentiate the identity
\[
\varphi(c^*(\theta),\theta)=0
\]
with respect to $\theta$. By the chain rule,
\[
\partial_c\varphi(c^*(\theta),\theta)\,\nabla_\theta c^*(\theta)
+
\nabla_\theta\varphi(c^*(\theta),\theta)
=
0.
\]
Evaluating at $\theta=\theta^*$ and solving for $\nabla_\theta c^*(\theta^*)$ yields
\[
\nabla_\theta c^*(\theta^*)
= -\frac{\nabla_\theta\varphi(c^*,\theta^*)}{\partial_c\varphi(c^*,\theta^*)},\]
which proves~\eqref{eq:gradient-cstar}. Substituting this into~\eqref{eq:V-general} gives~\eqref{eq:variance}. The explicit expressions follow from
\[
\partial_c\varphi(c^*,\theta^*)
=
\lambda_1\pi_1 \partial_c f_1(c^*;\theta_1^*)-\lambda_0\pi_0 \partial_c f_0(c^*;\theta_0^*)
\]
and
\[
\nabla_\theta\varphi(c^*,\theta^*)
=
\begin{pmatrix}
-\lambda_0\pi_0\,\nabla_{\theta_0}f_0(c^*;\theta_0^*)\\[0.4em]
\lambda_1\pi_1\,\nabla_{\theta_1}f_1(c^*;\theta_1^*)
\end{pmatrix}.
\]
\end{proof}

\end{document}
